\hypertarget{appendix-17fc1bc8b305}{}
\section*{Octal}
\section{Octal}\label{appendix-17fc1bc8b305-octal}

\textbf{Curriculum depth:} Introduced · \textbf{First seen:} C001 · \textbf{Latest development:} C001

\subsection{Definition and mechanism}\label{appendix-17fc1bc8b305-definition-and-mechanism}

Base 8 maps one digit to three bits; it survives in permissions and legacy literal syntax.

Octal digits run from 0 through 7. Group binary from the right in threes: \texttt{010\ 101\ 101₂\ =\ 255₈}. In C++, \texttt{0255} is octal 173; decimal \texttt{255} is a different value. Permission triples such as rwx map naturally onto three bits.

\subsection{Prerequisites}\label{appendix-17fc1bc8b305-prerequisites}

\begin{itemize}
\tightlist
\item
  \hyperref[appendix-83c3f28f7bda]{Positional notation} --- A base is the number of available digits.
\end{itemize}

\subsection{Used by}\label{appendix-17fc1bc8b305-used-by}

Connect later chunks here when they use this concept.

\subsection{Recall and next depth}\label{appendix-17fc1bc8b305-recall-and-next-depth}

Explain octal without the label, work the example above, then state the assumption whose removal would change the result.

\textbf{Pending:} Unix permission semantics.

\subsection{Source and navigation}\label{appendix-17fc1bc8b305-source-and-navigation}

\hyperref[01-foundations-i-ch001-from-a-signal-to-an-interpretation]{C001 teaching chapter} · \hyperref[appendix-fd090d982a15]{Part definitions and mechanisms} · \href{../02-Foundations-II/main.pdf\#nameddest=appendix-4ad47bdd269a}{Dependency Map (02-Foundations-II, appendix-4ad47bdd269a)}

\subsection{Why octal still matters}\label{appendix-17fc1bc8b305-why-octal-still-matters}

Octal groups binary digits in threes. Each octal digit therefore represents exactly three bits, so a binary field such as \texttt{111\ 101\ 010} maps directly to \texttt{752₈} without repeated division. Hexadecimal dominates modern memory dumps because bytes align naturally to two hex digits, but octal survives in domains whose historical or semantic fields are three-bit groups.

POSIX permission notation is the canonical example: read/write/execute are bit values \(4,2,1\), so permission \texttt{0754} means owner \texttt{111}, group \texttt{101}, other \texttt{100}. The leading zero in C/C++ source historically denotes an octal integer literal, which creates a real parsing hazard: \texttt{010} is eight, not ten. Modern code should prefer explicit intent and avoid visually ambiguous numeric literals.

A reconstruction test should convert binary↔octal by bit grouping, explain why each digit is 0--7, decode a permission triple, and identify the literal-syntax hazard. Octal is not a different storage representation in memory; it is a textual interpretation of the same integer bits.
